\textit{Demostración:}\\
Partimos de la definición de probabilidad condicional:

\[
\text{(1)} \qquad
p(Y = y \mid X = x)
=
\frac{p(X = x \cap Y = y)}{p(X = x)}
\]

\[
\text{(2)} \qquad
p(X = x \mid Y = y)
=
\frac{p(X = x \cap Y = y)}{p(Y = y)}
\]

Despejamos la probabilidad conjunta en (2):

\[
p(X = x \cap Y = y)
=
p(X = x \mid Y = y)\,p(Y = y)
\]

y sustituimos en (1):

\[
p(Y = y \mid X = x)
=
\frac{p(X = x \mid Y = y)\,p(Y = y)}{p(X = x)}
\]

Por teorema de probabilidad total tenemos que:

\[
p(X = x)
=
\sum_{y_i \in Y} p(X = x \cap Y = y_i)
=
\sum_{y_i \in Y} p(X = x \mid Y = y_i)\,p(Y = y_i)
\]

Sustituimos \(p(X = x)\):

\[
p(Y = y \mid X = x)
=
\frac{
p(X = x \mid Y = y)\,p(Y = y)
}{
\displaystyle \sum_{y_i \in Y} p(X = x \mid Y = y_i)\,p(Y = y_i)
}
\hspace{2cm} \blacksquare
\]

